% Nouveau EGA V (cote 25): archive cote 25, site page 51 (PDF page 52).
% typescript. Draft compared: none.
% AI-transcribed by Claude Opus 5.5 (claude-opus-5-5, adaptive thinking, effort high), Message Batches API, batch ega25b (msgbatch_01A7n5YuXubUPNK1zsepi8RG), 2026-11-13; not yet checked by eye.
\scanpage{25/51}
\var{\margin{\unsure{16}}}{\placedmargin{\unsure{16}}}{marginal note placed}{ega25_25_051-v0001}
\var{\margin{o--- c)}}{\placedmargin{o--- c)}}{marginal note placed}{ega25_25_051-v0002} \var{\inserted{Supposons cependant X lisse de dim 2, géom. intègre \unsure{propre} / k, et f une \uline{immersion}. Alors il résulte de 6.1. et 8.8. \var{\inserted{et 8.14.}}{et 8.14.}{insertion}{ega25_25_051-v0004} \unsure{que} si $Y^{\mathrm{sing}}_{\xi} \to \check{P}$ n'est pas gén. fini et inséparable, l'ens. des $\{\xi \in \check{P}$ tel que $Y_{\xi}$ soit \var{\deleted{\unsure{géom.}}}{}{deletion}{ega25_25_051-v0005} séparable sur $k(\xi)\}$ avec \uline{au plus 2} comp. irréd. géom. -- un compl. de codim $\geq 2$}}{Supposons cependant X lisse de dim 2, géom. intègre \unsure{propre} / k, et f une \uline{immersion}. Alors il résulte de 6.1. et 8.8. \var{\inserted{et 8.14.}}{et 8.14.}{insertion}{ega25_25_051-v0004} \unsure{que} si $Y^{\mathrm{sing}}_{\xi} \to \check{P}$ n'est pas gén. fini et inséparable, l'ens. des $\{\xi \in \check{P}$ tel que $Y_{\xi}$ soit \var{\deleted{\unsure{géom.}}}{}{deletion}{ega25_25_051-v0005} séparable sur $k(\xi)\}$ avec \uline{au plus 2} comp. irréd. géom. -- un compl. de codim $\geq 2$}{insertion}{ega25_25_051-v0003}
une quadrique de l'espace dual, relative à la forme duale ...).   Dans le
cas d'une surface non singulière \var{\xout{dx}}{}{typewriter strike-out}{ega25_25_051-v0006} d'un espace projectif, cette situa-%
tion doit cependant être considérée comme exceptionnelle, cf nº suivant.

\var{\margin{o---}}{\placedmargin{o---}}{marginal note placed}{ega25_25_051-v0007}
Nous allons examiner de plus près   maintenant le cas des surfaces
(le cas des courbes ne se pose évidemment pas du point de vue irréductibi-%
lité des sections hyperplanes \var{\xout{\unsure{t}}}{}{typewriter strike-out}{ega25_25_051-v0008} ).

(NB  Je m'aperçois avec effroi que la quadrique n'a pas droit, dans
l'emploi que j'ai fait du mot ``\var{\xout{f}}{}{typewriter strike-out}{ega25_25_051-v0009}réglé'', à être appelée réglée !  Cela doit
être en désaccord avec nos pères, et il faudra donc inventer un mot plus
adéquat pour la notion utilisée ici).

\uline{Proposition} 8.18.  Supposons \var{\xout{X géom intègre}}{}{typewriter strike-out}{ega25_25_051-v0010} k alg clos, X \var{\xout{intègre et}}{}{typewriter strike-out}{ega25_25_051-v0011}
\var{\xout{et}}{}{typewriter strike-out}{ega25_25_051-v0012} intègre (resp. intègre et normale)\var{\inserted{de dimension $d \geq 2$ et}}{de dimension $d \geq 2$ et}{insertion}{ega25_25_051-v0013}
propre sur k, soit T une partie finie fermée de X telle que  X-T soit
lisse  et f $|$X-T  soit non ramifié.  Pour que l'ensemble des \hand{$\xi \in$} $\check{P}$  tels
que  $Y_{\hand{\xi}}$  soit \var{\xout{\unsure{es}}}{}{typewriter strike-out}{ega25_25_051-v0014} géom irréd   (resp.  géom intègre)  de dimension  1
\unsure{fait}  un complémentaire de  codimension $\geq$ 2, il f et s que les deux
conditions suivantes soient satisfaites:

a)  Pour tout x\hand{$\in$}T, il existe une section hyperplane $Y_{\hand{\xi}}$   (\hand{$\xi \in$}P(k))
passant par x, de dimension \var{\over{1}{$d-1$}}{$d-1$}{writing over}{ega25_25_051-v0015}, et qui soit irréductible.

b) \var{\inserted{X'=}}{X'=}{insertion}{ega25_25_051-v0016} X-T est ``réglé'' (sic)  pour f, \uline{ou}  il existe  une section hyper
plane  $Y'_{\hand{\xi}}$  (\hand{$\xi \in$} $\check{P}$(k) )  de X',  \var{\xout{singulière}}{}{typewriter strike-out}{ega25_25_051-v0017} qui soit de dimension \var{\over{1}{d}}{d}{writing over}{ega25_25_051-v0018}\hand{-1},
singulière et  irréductible.

\var{\deleted{Soit  en effet  Z  la  réunion  de  g($\overline{\underline{Y}'^{\hand{\mathrm{sing}}}}$), \var{\xout{et}}{}{typewriter strike-out}{ega25_25_051-v0020} des hyperplans
de P  définis par les  f(x), x  T.  \var{\xout{et}}{}{typewriter strike-out}{ega25_25_051-v0021}  Alors pour   P - Z,  Y   est
\uline{lisse} de dimension 1, \var{\xout{\illegible{8}}}{}{typewriter strike-out}{ega25_25_051-v0022}
\var{\xout{\illegible{10} au même \illegible{2} est lisse au dess}}{}{typewriter strike-out}{ega25_25_051-v0023}
\var{\inserted{géom}}{géom}{insertion}{ega25_25_051-v0024}
\var{\xout{\unsure{des points de}}}{}{typewriter strike-out}{ega25_25_051-v0025}  et comme  il est connexe par 6.1. , il est géom irréd.}}{}{deletion}{ega25_25_051-v0019}

\mspage{51}
